\documentclass[11pt,a4paper]{article}
\usepackage[margin=2cm]{geometry}
\usepackage[T1]{fontenc}
\usepackage{mathptmx}
\usepackage{amsmath,amssymb}
\usepackage{xcolor}
\usepackage{enumitem}
\usepackage{array}
\usepackage{float}
\usepackage{colortbl}
\usepackage{needspace}
\usepackage[hidelinks]{hyperref}
\definecolor{navy}{HTML}{1B3755}
\definecolor{headc}{HTML}{E6ECF5}
\definecolor{okc}{HTML}{DFF3DF}
\definecolor{noc}{HTML}{FBE0E0}
\definecolor{warnc}{HTML}{FFF8E8}
\setlength{\parindent}{0pt}
\setlength{\parskip}{5pt}
\renewcommand{\arraystretch}{1.25}
\newcommand{\sect}[1]{\needspace{8\baselineskip}\bigskip{\Large\color{navy}\bfseries #1}\par\smallskip{\color{navy}\hrule height 0.8pt}\smallskip}
\newcommand{\sub}[1]{\needspace{5\baselineskip}\medskip{\large\color{navy}\bfseries #1}\par}
\newcommand{\hd}{\rowcolor{headc}}
\newcommand{\cmd}[1]{\par\smallskip\noindent\colorbox{headc}{\parbox{\dimexpr\linewidth-2\fboxsep}{\ttfamily\small #1}}\par\smallskip}
\newcommand{\note}[1]{\par\smallskip\noindent\fcolorbox{orange}{warnc}{\parbox{\dimexpr\linewidth-2\fboxsep-2\fboxrule}{#1}}\par\smallskip}
\newcommand{\say}[1]{\par\noindent\textit{Say: ``#1''}\par}

\begin{document}
\begin{center}
{\LARGE\bfseries\color{navy} Running the Model in Front of the Examiner}\\[4pt]
{\large A step-by-step guide for the defence}\\[2pt]
IEEE 13-node feeder, DIgSILENT PowerFactory 2021 --- Jhala Nath Kafle (081MSPSE009)
\end{center}

\note{\textbf{The short version.} Open PowerFactory, activate the project \textbf{IEEE13 DSDR Fuse Coordination} and
the study case \textbf{Study with Substation Transformer}. Then show three things by hand: (1) the load flow
with the DG out and in --- the current through R2 reverses; (2) a short circuit at node 633 --- the fuse
carries more current than the recloser; (3) the tripping times of the relays and the fuse. That takes about
five minutes and proves the model is real. The scripts are the backup, not the main show.}

\sect{1. Before the defence (do this the day before)}
\begin{enumerate}[leftmargin=*,itemsep=2pt]
\item \textbf{Use the same laptop} you built the model on. The project is stored in PowerFactory's own database
on this machine, not in the project folder.
\item \textbf{Check the licence.} Open PowerFactory once where you will present. If the licence needs a
network or a dongle, make sure it works there.
\item \textbf{Rehearse every step of Section 2 once}, and write down the numbers you see next to the numbers in
Section 5. The menu names below are those of PowerFactory 2021; if a button is somewhere else on your
screen, note where.
\item \textbf{Leave the model in the DSDR state.} It is left that way after the last run of the study.
\item \textbf{Keep a backup:} in PowerFactory, right-click the project \,$\rightarrow$\, \emph{Export}, and save
the \texttt{.pfd} file on a USB stick.
\item \textbf{Have these open and minimised:} PowerFactory with the project active; the report PDF; and
the figures folder, \texttt{results\textbackslash figures}, inside \texttt{DSDR\_Study}.
\end{enumerate}

\sect{2. Showing the model by hand in PowerFactory (recommended)}

\sub{Step 1 --- Open and activate (30 seconds)}
\begin{enumerate}[leftmargin=*,itemsep=2pt]
\item Start \textbf{PowerFactory 2021}.
\item Open the Data Manager, find the project \textbf{IEEE13 DSDR Fuse Coordination} under your user,
right-click \,$\rightarrow$\, \textbf{Activate}.
\item In the study-case list at the top, choose \textbf{Study with Substation Transformer}.
\item The single-line diagram of the IEEE 13-node feeder appears.
\end{enumerate}
\say{This is the IEEE 13-node feeder at 4.16 kV, fed from a 115 kV grid through a 5 MVA transformer. R1 is
here at the feeder head, R2 is on line 632 to 671, the fuses are on the laterals, and the DG is at node 692.}

\sub{Step 2 --- Load flow without the DG (1 minute)}
\begin{enumerate}[leftmargin=*,itemsep=2pt]
\item Right-click the \textbf{Synchronous Machine} at node 692 \,$\rightarrow$\, set it \textbf{Out of Service}.
\item Click \textbf{Calculate Load Flow}. Select \emph{AC Load Flow, unbalanced, 3-phase (ABC)} and
\textbf{Execute}.
\item Point at the result boxes of line 650--632 and of line 632--671.
\end{enumerate}
\say{Without DG the feeder head carries about 588 amperes and the branch of R2 about 470 amperes, flowing
from the grid towards the loads.}

\sub{Step 3 --- Load flow with the DG (1 minute)}
\begin{enumerate}[leftmargin=*,itemsep=2pt]
\item Put the Synchronous Machine back \textbf{In Service}.
\item \textbf{Calculate Load Flow} again.
\item Point at line 632--671: the current is now about 252 amperes and the power flows \emph{towards 632}.
\end{enumerate}
\say{With the DG the current through R2 falls to about 252 amperes and its direction reverses. This is why
R2 needs a second, reverse setting: 1.25 times 252 gives the reverse pickup of 315 amperes.}

\sub{Step 4 --- Short circuit (1 to 2 minutes)}
\begin{enumerate}[leftmargin=*,itemsep=2pt]
\item Right-click the busbar \textbf{633} \,$\rightarrow$\, \textbf{Calculate} \,$\rightarrow$\,
\textbf{Short-Circuit}.
\item Method: \textbf{complete}. Fault type: \textbf{3-phase short-circuit}. Fault impedance 0.
\textbf{Execute}.
\item Point at the currents in the fuse F633, in R1 (line 650--632) and in R2 (line 632--671).
\end{enumerate}
\say{For a three-phase fault at 633 with the DG, the fuse carries about 5.4 kiloamperes, but R1 carries
only about 4.0 and R2 about 1.5 in reverse. The fuse sees the grid and the DG current together; each
recloser sees only its share. That is the whole problem.}

\sub{Step 5 --- Tripping times (1 minute)}
With the short circuit of Step 4 still calculated:
\begin{enumerate}[leftmargin=*,itemsep=2pt]
\item Open the relay \textbf{R1 Fast} (double-click it in the cubicle, or find it in the Data Manager) and
look at its \textbf{tripping time}. Do the same for \textbf{R2 Rev Fast} and for the fuse \textbf{F633}.
\item Or open a \textbf{time-overcurrent plot} that contains these devices: the vertical line is the fault
current, and the plot shows where it meets each curve.
\end{enumerate}
\say{R2 trips on its reverse group in about 0.06 seconds and R1 in about 0.11 seconds. The fuse would start
to melt at about 0.15 seconds. So both reclosers trip first and the fuse is saved.}

\sub{Step 6 --- If asked for the EMT simulation}
The EMT run needs the scripted reclosing events, so it is quicker to show the saved result: Figure 5.14 of
the report, or the picture file
\cmd{DSDR\_Study\textbackslash results\textbackslash figures\textbackslash Fig10\_EMT\_LL\_684.png}
If the examiner insists on a live run, use Section 3, command (d); it runs the simulation six times and
takes a few minutes.

\sect{3. Running the scripts (backup, or if asked ``show me the code running'')}
\note{\textbf{Close PowerFactory first.} The scripts start their own PowerFactory in the background, and that
cannot start while the PowerFactory window is open.}

Open a Command Prompt and go to the study folder:
\cmd{cd /d "C:\textbackslash Users\textbackslash Rabin\textbackslash Desktop\textbackslash Digsilent\_project - TWIST AND TURN\textbackslash DSDR\_Study"}
\cmd{set PY="C:\textbackslash Users\textbackslash Rabin\textbackslash AppData\textbackslash Local\textbackslash Programs\textbackslash Python\textbackslash Python39\textbackslash python.exe"}

\begin{table}[H]\centering\small
\begin{tabular}{|c|p{6.2cm}|p{8cm}|}\hline
\hd & \textbf{Command} & \textbf{What it does and what appears} \\ \hline
(a) & \texttt{\%PY\% step2\_studies.py} & Load flows and all short circuits in PowerFactory. Prints the
currents through R1 and R2 for each fault. \\ \hline
(b) & \texttt{\%PY\% step3\_design\_and\_evaluate.py} & No PowerFactory needed; a few seconds. Prints Table II,
the settings, and the cell counts: 35 and 39 without DG, 24 with DG, 38 with the DSDR. \\ \hline
(c) & \texttt{\%PY\% step4\_apply\_settings.py dsdr} & Writes the DSDR settings into the relays and fuses and
compares PowerFactory's tripping times with the calculated ones. \\ \hline
(d) & \texttt{\%PY\% step5\_fig10\_emt.py} & The EMT simulation, DG out and DG in. Prints the event times and
saves the figure. \\ \hline
\end{tabular}
\end{table}

\textbf{The safest script to run live is (b).} It needs no PowerFactory, finishes in seconds and prints the
headline results. The full output is also saved in \texttt{results\textbackslash summary.txt}, which you can
simply open.

\note{\textbf{Do not run \texttt{run\_all.py} or \texttt{step1\_build\_model.py} at the defence.} They rebuild
the model and all outputs from scratch, which takes several minutes and is not needed to show a result.}

\sect{4. A third way: the script inside PowerFactory}
The project contains a Python script object named \textbf{Final(1)}. With the project active, find it in the
Data Manager (under the study case or the project's \emph{Scripts}), right-click \,$\rightarrow$\,
\textbf{Execute}. It runs the load flows and the short-circuit database and prints a report in the
\textbf{Output Window}, with the current through R1 and R2 and its direction, DG out and DG in. Rehearse
it once; it takes a few minutes.

\sect{5. Numbers you should see}
\begin{table}[H]\centering\small
\begin{tabular}{|p{6.6cm}|c|c|}\hline
\hd \textbf{Load flow} & \textbf{DG out} & \textbf{DG in} \\ \hline
Feeder head, line 650--632 (R1) & 587.7 A & 284.9 A \\ \hline
Line 632--671 (R2) & 470.0 A, forward & 252.2 A, \textbf{reverse} \\ \hline
\end{tabular}
\end{table}
\begin{table}[H]\centering\small
\begin{tabular}{|p{6.6cm}|c|c|}\hline
\hd \textbf{Bolted three-phase fault} & \textbf{DG out} & \textbf{DG in} \\ \hline
Fault current at 632 & 4.74 kA & 6.53 kA \\ \hline
Fault current at 633 & 4.10 kA & 5.43 kA \\ \hline
\end{tabular}
\end{table}
\begin{table}[H]\centering\small
\begin{tabular}{|p{4.2cm}|c|c|c|}\hline
\hd \textbf{Fault at 633, DG in, DSDR} & \textbf{Current} & \textbf{Device} & \textbf{Time} \\ \hline
R2, reverse group & 1500 A & R2 Rev Fast & 0.060 s \\ \hline
R1 & 3970 A & R1 Fast & 0.106 s \\ \hline
Fuse F633 (400E) & 5426 A & F633 & melts at 0.149 s \\ \hline
\end{tabular}
\end{table}
Small differences in the last digit are normal. If a number is very different, check the two things in
Section 6 first.

\sect{6. If something goes wrong}
\begin{table}[H]\centering\small
\begin{tabular}{|p{5.4cm}|p{9.9cm}|}\hline
\hd \textbf{What you see} & \textbf{What to do} \\ \hline
The numbers are much too high (for example 6 to 8 kA at 632 without DG) & The wrong study case is active.
Choose \textbf{Study with Substation Transformer}. \\ \hline
R2 current is 470 A although you expect 252 A & The DG is out of service. Put the Synchronous Machine in
service and run the load flow again. \\ \hline
Load flow does not converge & Press \emph{Reset Calculation}, check that nothing was switched off by mistake,
run again. \\ \hline
``Short-circuit calculation not possible'' & You chose a phase the node does not have (645 and 646 have b--c
only, 684 has a--c, 611 has c, 652 has a). Choose node 633, 632 or 671. \\ \hline
A script says PowerFactory could not start & PowerFactory's window is still open. Close it and run the
command again. \\ \hline
A script stops with an error & Do not debug in front of the panel. Open \texttt{results\textbackslash summary.txt}
and the saved figures, and say the results were produced by the same script. \\ \hline
PowerFactory will not open at all & Show the report figures and \texttt{results\textbackslash summary.txt}; offer
to run it afterwards. \\ \hline
\end{tabular}
\end{table}

\sect{7. What the examiner may ask while it runs}
\textbf{Which calculation method is this?} Unbalanced, three-phase load flow; short circuit by the
\emph{complete} method, which starts from the load-flow state.

\textbf{Why this study case and not the other one?} This one includes the 115/4.16 kV, 5 MVA substation
transformer. With it the fault levels match the IEEE short-circuit benchmark within about 2 per cent. The
stiff-grid case gives fault levels that are too high.

\textbf{Where are the two settings of R2?} The CDG34 relay of the library is not directional. So R2 has
separate relay units: \emph{R2 Fast} and \emph{R2 Delayed} for the forward group, \emph{R2 Rev Fast} and
\emph{R2 Rev Delayed} for the reverse group. The direction is assigned from the fault location. I state
this as a limitation.

\textbf{What do the Python scripts do that PowerFactory does not?} PowerFactory does the load flow, the short
circuits and the EMT simulation. The scripts run those studies for all nodes and fault types, apply the
method's equations to get the settings, and classify the 39 cells. Doing 273 short circuits by hand would
not be practical.

\textbf{Can you change the DG size and show the effect?} Open the Synchronous Machine, change its active
power, run the load flow: the current through R2 changes. The full study at seven penetration levels is in
Section 5.16 of the report.

\end{document}
